% AUTO-GENERATED by paper/gen_data.py from docs/data/levels.json — do not edit.
\begin{longtable}{@{}c l p{8cm}@{}}
\toprule
Lvl & Notation & How much is that \\
\midrule
\endhead
a & $a$ & Overtakes a googol and every physically meaningful count. \\
b & $a^{\star 1}$ & Past a googolplex — its exponent alone is astronomical. \\
c & $a^{\star 2}$ & Its number of digits has more digits than the universe has atoms. \\
d & $a^{\star 3}$ & Take the logarithm three times over — still beyond all physics. \\
e & $a^{\star 4}$ & Raise a googolplex to a googolplex; you are not even close. \\
f & $a^{\star 5}$ & More than Shannon's chess games raised to the positions of Go. \\
g & $a^{\star 6}$ & A busy-beaver machine would halt before printing the first exponent. \\
h & $a^{\star 7}$ & Write Avogadro's number in every Planck volume of the cosmos — not close. \\
i & $a^{\star 8}$ & Every digit a universe, every universe's atoms digits — still short. \\
j & $a^{\star 9}$ & The number of its digits is itself beyond counting. \\
k & $a^{\star 10}$ & Deep in tetration: self-powers stacked past every named number. \\
l & $a^{\star 11}$ & The Ackermann function would treat this as an early, warm-up input. \\
m & $a^{\star 12}$ & Beyond Skewes's number, once called 'the largest in mathematics'. \\
n & $a^{\star 13}$ & Name a new number every Planck time since the Big Bang — never reach it. \\
o & $a^{\star 14}$ & A googolplexian is a rounding error down at its base. \\
p & $a^{\star 15}$ & Its logarithm needs its own logarithm, which needs its own… \\
q & $a^{\star 16}$ & Taller than any tower of tens you could build atom by atom. \\
r & $a^{\star 17}$ & The exponent of the exponent of the exponent is already unphysical. \\
s & $a^{\star 18}$ & Past every 'illion' and every number humans have ever named. \\
t & $a^{\star 19}$ & You'd exhaust factorials of factorials before catching up. \\
u & $a^{\star 20}$ & Twenty rounds of 'raise to itself', and you have barely arrived. \\
v & $a^{\star 21}$ & The power tower is now so tall the tens blur into a single line. \\
w & $a^{\star 22}$ & And yet — a humbling truth — infinitely short of Graham's number. \\
x & $a^{\star 23}$ & Infinitely short of TREE(3), which dwarfs Graham's number. \\
y & $a^{\star 24}$ & The last letter — the alphabet ran out before the number did. \\
z & $a^{\star 25}$ & The summit — yet still a finite number you could, in principle, write. \\
$\ccoma$ & $a^{\star 26}$ & The summit — yet still a finite number you could, in principle, write. \\
\bottomrule
\end{longtable}
